
%\input{derivmacros}
%% Letters %%%%%%%%%%%%%%%%%%%%%%%%%%

\def\tA{\tilde{A}}
\def\tB{\tilde{B}}
\def\tC{\tilde{C}}
\def\tP{\tilde{P}}
\def\tQ{\tilde{Q}}
\def\tR{\tilde{R}}
\def\tS{\tilde{S}}
\def\tT{\tilde{T}}
\def\tU{\tilde{U}}
\def\tW{\tilde{W}}
\def\tDelta{\tilde{\Delta}}
\def\tGamma{\tilde{\Gamma}}
\def\tPi{\tilde{\Pi}}
\def\tpi{\tilde{\pi}}
\def\tsigma{\tilde{\sigma}}
\def\ts{\tilde{s}}

\def\ha{\hat{a}}

\def\hA{\hat{A}}
\def\hB{\hat{B}}
\def\hP{\hat{P}}
\def\hQ{\hat{Q}}
\def\hR{\hat{R}}
\def\hS{\hat{S}}
\def\hT{\hat{T}}
\def\hW{\hat{W}}
\def\hDelta{\hat{\Delta}}
\def\hGamma{\hat{\Gamma}}
\def\hrho{\hat{\rho}}
\def\hphi{\hat{\phi}}
\def\hpsi{\hat{\psi}}
\def\vphi{\check{\phi}}

\def\cA{{\mathcal A}}
\def\cB{{\mathcal B}}
\def\cC{{\mathcal C}}
\def\cD{{\mathcal D}}
\def\cE{{\mathcal E}}
\def\cF{{\mathcal F}}
\def\cG{{\mathcal G}}
\def\cH{{\mathcal H}}
\def\cI{{\mathcal I}}
\def\cJ{{\mathcal J}}
\def\cK{{\mathcal K}}
\def\cL{{\mathcal L}}
\def\cM{{\mathcal M}}
\def\cN{{\mathcal N}}
\def\cO{{\mathcal O}}
\def\cP{{\mathcal P}}
\def\cQ{{\mathcal Q}}
\def\cR{{\mathcal R}}
\def\cS{{\mathcal S}}
\def\cT{{\mathcal T}}
\def\cU{{\mathcal U}}
\def\cV{{\mathcal V}}
\def\cW{{\mathcal W}}
\def\cX{{\mathcal X}}
\def\cY{{\mathcal Y}}
\def\cZ{{\mathcal Z}}

\newcommand{\fA}{{\mathfrak A}}
\newcommand{\fB}{{\mathfrak B}}
\newcommand{\fC}{{\mathfrak C}}
\newcommand{\fD}{{\mathfrak D}}
\newcommand{\fE}{{\mathfrak E}}
\newcommand{\fF}{{\mathfrak F}}
\newcommand{\fG}{{\mathfrak G}}
\newcommand{\fH}{{\mathfrak H}}
\newcommand{\fI}{{\mathfrak I}}
\newcommand{\fJ}{{\mathfrak J}}
\newcommand{\fK}{{\mathfrak K}}
\newcommand{\fL}{{\mathfrak L}}
\newcommand{\fM}{{\mathfrak M}}
\newcommand{\fN}{{\mathfrak N}}
\newcommand{\fO}{{\mathfrak O}}
\newcommand{\fP}{{\mathfrak P}}
\newcommand{\fQ}{{\mathfrak Q}}
\newcommand{\fR}{{\mathfrak R}}
\newcommand{\fS}{{\mathfrak S}}
\newcommand{\fT}{{\mathfrak T}}
\newcommand{\fU}{{\mathfrak U}}
\newcommand{\fV}{{\mathfrak V}}
\newcommand{\fW}{{\mathfrak W}}
\newcommand{\fX}{{\mathfrak X}}
\newcommand{\fY}{{\mathfrak Y}}
\newcommand{\fZ}{{\mathfrak Z}}
\newcommand{\fa}{{\mathfrak a}}
\newcommand{\fb}{{\mathfrak b}}
\newcommand{\fc}{{\mathfrak c}}
\newcommand{\fd}{{\mathfrak d}}
\newcommand{\fe}{{\mathfrak e}}
\newcommand{\ff}{{\mathfrak f}}
\newcommand{\fg}{{\mathfrak g}}
\newcommand{\fh}{{\mathfrak h}}
%\newcommand{\fi}{{\mathfrak i}}
\newcommand{\fj}{{\mathfrak j}}
\newcommand{\fk}{{\mathfrak k}}
\newcommand{\fl}{{\mathfrak l}}
\newcommand{\fm}{{\mathfrak m}}
\newcommand{\fn}{{\mathfrak n}}
\newcommand{\fo}{{\mathfrak o}}
\newcommand{\fp}{{\mathfrak p}}
\newcommand{\fq}{{\mathfrak q}}
\newcommand{\fr}{{\mathfrak r}}
\newcommand{\fs}{{\mathfrak s}}
\newcommand{\ft}{{\mathfrak t}}
\newcommand{\fu}{{\mathfrak u}}
\newcommand{\fv}{{\mathfrak v}}
\newcommand{\fw}{{\mathfrak w}}
\newcommand{\fx}{{\mathfrak x}}
\newcommand{\fy}{{\mathfrak y}}
\newcommand{\fz}{{\mathfrak z}}

\def\Bool{\mathbb{B}}
\def\Nat{\mathbb{N}}

%%%%%%%%%%%%%%%%%%%%%%%%%%%%%%%%%%%%%%%%
% names of systems                    %%

\def\SKS{\mathsf{SKS}}
\def\eSKS{\mathsf{eSKS}}
\def\sSKS{\mathsf{sSKS}}
\def\SKSd{\mathsf{SKS}\mathord{\downarrow}}
\def\SKSu{\mathsf{SKS}\mathord{\uparrow}}
\def\KS{\mathsf{KS}}
\def\SKSg{\mathsf{SKSg}}
\def\SKSgd{\mathsf{SKSg}\mathord{\downarrow}}
\def\SKSgu{\mathsf{SKSg}\mathord{\uparrow}}
\def\KSg{\mathsf{KSg}}
\def\K{\mathsf{K}}
\def\MLL{\mathsf{MLL}}
\def\MALL{\mathsf{MALL}}


%%%%%%%%%%%%%%%%%%%%%%%%%%%%%%%%%%%%%%%%

\def\set#1{\{#1\}}
\def\bigset#1{\big\{#1\big\}}
\def\dset#1{\left\{#1\right\}}
\def\mset#1{\{#1\}}
\def\cons#1{\{#1\}}
\def\conhole      {\cons{\enspace}}%

\def\disunion{\mathbin{\uplus}}
\def\directsum{\mathbin{\oplus}}

\def\fcomp{\mathbin{\circ}}

\def\scope{\mathrel{.}}

\def\tuple#1{\langle#1\rangle}

\def\pair#1#2{(#1,#2)}
\def\npair#1#2{#1,#2}

%% relations %%%%%%%%%%%%%%%%%%%%%%%%%%%%%%%%%

\def\grammareq {\mathrel{\raise.4pt\hbox{::}{=}}}%

%%%%%%%%%%%%%%%%%%%%%%%%%%%%%%%%%%%%%%
%% category theory

\def\id{{\mathrm{1}}}
\def\idh{{\hat{\id}}}
\def\idch{{\check{\id}}}
%\def\mix{{\mathrm{mix}}}
%\def\twist{{\mathrm{twist}}}

\def\codiag{\nabla}
\def\diag{\Delta}
\def\proj{\Pi}
\def\coproj{\amalg}

\newcommand{\sto}[1][]{\stackrel{#1}{\to}}
\newcommand{\lsto}[1][]{\stackrel{#1}{\longrightarrow}}
\newcommand{\lsfrom}[1][]{\stackrel{#1}{\longleftarrow}}

\newcommand{\dotto}[1][]{\mathrel{\!\xy\ar@{.>}^-{#1}(5,0)\endxy\!}}
\newcommand{\solto}[1][]{\mathrel{\!\xy\ar@{->}^-{#1}(5,0)\endxy\!}}
\newcommand{\longsolto}[1][]{\mathrel{\!\xy\ar@{->}^-{#1}(11,0)\endxy\!}}
\newcommand{\longdotto}[1][]{\mathrel{\!\xy\ar@{.>}^-{#1}(11,0)\endxy\!}}
\newcommand{\xldotto}[2][]{\mathrel{\!\xy\ar@{.>}^-{#1}(#2,0)\endxy\!}}

\newcommand{\downmap}[1]{\begin{array}{c}#1\end{array}}
\newcommand{\vcdownmap}[1]{\vcenter{\begin{array}{c}#1\end{array}}}
\newcommand{\ldownto}[1][]{
  \\
  \llap{$\scriptstyle #1\;$}\clapm{\Big\downarrow}
  \\%[.2pt]
  }
\newcommand{\rdownto}[1][]{
  \\
  \clapm{\Big\downarrow}\rlap{$\scriptstyle\;#1$}
  \\%[.2pt]
  }

\def\biproj{p}
\def\bicoproj{i}
\def\bidiag{\blacktriangle}
\def\bicodiag{\blacktriangledown}
%\def\bidsum{\oplus}

%\def\runit{\varrho}
%\def\runitast{\varrho^\ast}

\def\Hom{}
\renewcommand{\Hom}[1][]{\mathrm{Hom}_{#1}}
%\def\Pom{}
%\renewcommand{\Pom}[1][]{\mathrm{Pom}_{#1}}
%\def\Mop{}
%\renewcommand{\Mop}[1][]{\mathrm{Mop}_{#1}}
\def\Obj{{\mathrm{Obj}}}
\def\isom{\cong}

\def\name#1{\hat{#1}}
\def\coname#1{\check{#1}}
\def\widename#1{\widehat{#1}}
\def\wideconame#1{\widecheck{#1}}

\def\Set{{\mathbf{Set}}}

%%%%%%%%%%%%%%%%%%%%%%%%%%%%%%%%%%%%%%%%%%%
% proof theory

\def\sqn  #1{{\;\turnstile #1\;}}%
\def\ssqn#1#2{{\;#1\turnstile #2\;}}%

\def\cutr{\mathsf{cut}}
\def\idr{\mathsf{id}}
\def\mixr{\mathsf{mix}}
\def\mixzr{\mathsf{mix_0}}
\def\weakr{\mathsf{weak}}
\def\conr{\mathsf{cont}}
\def\exr{\mathsf{exch}}
\def\cyclr{\mathsf{cycl}}
\def\rr{\mathsf{r}}
\def\swir{\mathsf{s}}
\def\medr{\mathsf{m}}

\def\MLL{\mathsf{MLL}}
\def\LK{\mathsf{LK}}
\def\Gthree{\mathsf{G3}}
\def\LC{\mathsf{LC}}
\def\LLP{\mathsf{LLP}}

% rule names

\newcommand{\ird}{\mathsf{i}\mathord{\downarrow}}
\newcommand{\iru}{\mathsf{i}\mathord{\uparrow}}
\newcommand{\atird}{\mathsf{ai}\mathord{\downarrow}}
\newcommand{\atiru}{\mathsf{ai}\mathord{\uparrow}}
\newcommand{\seqrd}{\mathsf{q}\mathord{\downarrow}}
\newcommand{\seqru}{\mathsf{q}\mathord{\uparrow}}
\newcommand{\conrd}{\mathsf{c}\mathord{\downarrow}}
\newcommand{\conru}{\mathsf{c}\mathord{\uparrow}}
\newcommand{\weakrd}{\mathsf{w}\mathord{\downarrow}}
\newcommand{\weakru}{\mathsf{w}\mathord{\uparrow}}
%\renewcommand{\swir}[1][]{\mathsf{s}_{#1}}
\newcommand{\assrd}[1][]{\alpha\mathord{\downarrow}_{#1}}
\newcommand{\assru}[1][]{\alpha\mathord{\uparrow}_{#1}}
\newcommand{\sassrd}[1][]{\alpha^\triangleleft\mathord{\downarrow}_{#1}}
\newcommand{\sassru}[1][]{\alpha^\triangleleft\mathord{\uparrow}_{#1}}
\newcommand{\comrd}[1][]{\sigma\mathord{\downarrow}_{#1}}
\newcommand{\comru}[1][]{\sigma\mathord{\uparrow}_{#1}}
\newcommand{\assoc}[1][]{\alpha_{#1}}
\newcommand{\twist}[1][]{\sigma_{#1}}

\def\ltensl{\mathord{\ltens}\mathsf{L}}
\def\ltensr{\mathord{\ltens}\mathsf{R}}
\def\lparl{\mathord{\lpar}\mathsf{L}}
\def\lparr{\mathord{\lpar}\mathsf{R}}
\def\limpl{\mathord{\limp}\mathsf{L}}
\def\limpr{\mathord{\limp}\mathsf{R}}
\def\lnegl{\mathord{\cdot^\lbot}\mathsf{L}}
\def\lnegr{\mathord{\cdot^\lbot}\mathsf{R}}
\def\exrl{\mathsf{exchL}}
\def\exrr{\mathsf{exchR}}
\def\markl{\mathord{\mathsf{:}}\mathsf{L}}
\def\markr{\mathord{\mathsf{:}}\mathsf{R}}

\def\poll{\bullet}
\def\polr{\circ}
\def\markpl{^\poll}
\def\markpr{^\polr}
\def\lnegpl{^{\lbot\poll}}
\def\lnegpr{^{\lbot\polr}}


\def\wlpar{\;\lpar\;}
\def\wltens{\;\ltens\;}

\def\un{\circ}


%%%%%%%%%%%%%%%%%%%%%%%%%%%%%%%%%%%%%%%%%%%
%% classical logic %%%%%%%%%%%%%%%%%%%%%%%%

\def\wvskipvi{\mathchoice{\hskip-2ex}{\hskip-1.5ex}{\hskip-1ex}{\hskip-.8ex}}
\def\wvskipx{\mathchoice{\hskip-2.5ex}{\hskip-2ex}{\hskip-1.5ex}{\hskip-1ex}}
\def\wvskipxii{\mathchoice{\hskip-2.9ex}{\hskip-2.25ex}{\hskip-2ex}{\hskip-1.5ex}}
\let\wvskip\wvskipx
\def\dbigand#1{\displaystyle\mathop{\bigwedge\wvskip\bigwedge}_{#1}}
\def\dbigor#1{\displaystyle\mathop{\bigvee\wvskip\bigvee}_{#1}}
\def\bigand#1{\mathop{\bigwedge\wvskip\bigwedge}_{\smash{#1}}}
\def\bigor#1{\mathop{\bigvee\wvskip\bigvee}_{\smash{#1}}}
\def\ctrue{{\mathbf t}}
\def\cfalse{{\mathbf f}}
\def\cand{\wedge}
\def\cor{\vee}
\def\candor{\mathchoice%
  {\mathbin{\rlap{$\vee$}\wedge}}
  {\mathbin{\rlap{$\vee$}\wedge}}
  {\mathbin{\rlap{$\scriptstyle\vee$}\wedge}}
  {\mathbin{\rlap{$\scriptscriptstyle\vee$}\wedge}}
}
\def\cplus{\oplus}
%\def\ccut{\varobar}
\def\ccon{\blacktriangledown}
\def\cneg#1{\bar{#1}}
\def\widecneg#1{\overline{#1}}
\def\wcneg#1{\overline{#1}}
\def\cimp{\mathop{\Rightarrow}}
\def\ciff{\mathop{\Leftrightarrow}}
\def\fneg{\widecneg{(-)}}

\newbox\cutbox
\newdimen\cutwd
\newdimen\cutht
\newdimen\cutdp
\def\ccut{%
  \setbox\cutbox\hbox{$\lozenge$}
  \cutwd=\wd\cutbox
  \cutht=\ht\cutbox
  \cutdp=\dp\cutbox
  \setbox\cutbox\hbox to\cutwd{\hss\vrule width.3pt height\cutht depth\cutdp\hss}
  %%%%%\mathbin{\rlap{$\lozenge$}\copy\cutbox}}
  \mathbin{\lozenge\hskip-\cutwd\copy\cutbox}}
\def\scriptcut{%
  \setbox\cutbox\hbox{$\scriptstyle\lozenge$}
  \cutwd=\wd\cutbox
  \cutht=\ht\cutbox
  \cutdp=\dp\cutbox
  \setbox\cutbox\hbox to\cutwd{\hss\vrule width.3pt height\cutht depth\cutdp\hss}
  \mathord{\lozenge\hskip-\cutwd\copy\cutbox}}

\def\vccut{%
  \setbox\cutbox\hbox{$\lozenge$}
  \cutwd=\wd\cutbox
  \cutht=\ht\cutbox
  \cutdp=\dp\cutbox
  \setbox\cutbox\hbox to\cutwd{\hss\hskip.3pt\vrule width.3pt height\cutht depth\cutdp\hss}
  %%%%%\mathbin{\rlap{$\lozenge$}\copy\cutbox}}
  \mathbin{\lozenge\hskip-\cutwd\copy\cutbox}}


%% CoS %%%%%%%%%%%%%%%%%%%%%%%%%%%%%%%%%%%%

\def\conldel {\{}%
\def\conrdel {\}}%
\def\lrgldel {\mathchoice{(}{(}{\langle}{\langle}}%
\def\lrgrdel {\mathchoice{)}{)}{\rangle}{\rangle}}%
\def\aprldel {\mathchoice
    {\mathopen {\setbox0=\hbox{$\displaystyle     \lrgldel$}\hbox to\wd0
                         {\hfil$\displaystyle     (       $\hfil}}}%
    {\mathopen {\setbox0=\hbox{$\textstyle        \lrgldel$}\hbox to\wd0
                         {\hfil$\textstyle        (        $\hfil}}}%
    {\mathopen {\setbox0=\hbox{$\scriptstyle      \lrgldel$}\hbox to\wd0
                         {\hfil$\scriptstyle      (        $\hfil}}}%
    {\mathopen {\setbox0=\hbox{$\scriptscriptstyle\lrgldel$}\hbox to\wd0
                         {\hfil$\scriptscriptstyle(        $\hfil}}}}%
\def\aprrdel {\mathchoice
    {\mathclose{\setbox0=\hbox{$\displaystyle     \lrgrdel$}\hbox to\wd0
                         {\hfil$\displaystyle     )       $\hfil}}}%
    {\mathclose{\setbox0=\hbox{$\textstyle        \lrgrdel$}\hbox to\wd0
                         {\hfil$\textstyle        )        $\hfil}}}%
    {\mathclose{\setbox0=\hbox{$\scriptstyle      \lrgrdel$}\hbox to\wd0
                         {\hfil$\scriptstyle      )        $\hfil}}}%
    {\mathclose{\setbox0=\hbox{$\scriptscriptstyle\lrgrdel$}\hbox to\wd0
                         {\hfil$\scriptscriptstyle)        $\hfil}}}}%
\def\seqldel {\mathchoice
    {\mathopen {\setbox0=\hbox{$\displaystyle     \lrgldel$}\hbox to\wd0
                         {\hfil$\displaystyle     \langle  $\hfil}}}%
    {\mathopen {\setbox0=\hbox{$\textstyle        \lrgldel$}\hbox to\wd0
                         {\hfil$\textstyle        \langle  $\hfil}}}%
    {\mathopen {\setbox0=\hbox{$\scriptstyle      \lrgldel$}\hbox to\wd0
                         {\hfil$\scriptstyle      \langle  $\hfil}}}%
    {\mathopen {\setbox0=\hbox{$\scriptscriptstyle\lrgldel$}\hbox to\wd0
                         {\hfil$\scriptscriptstyle\langle  $\hfil}}}}%
\def\seqrdel {\mathchoice
    {\mathclose{\setbox0=\hbox{$\displaystyle     \lrgrdel$}\hbox to\wd0
                         {\hfil$\displaystyle     \rangle  $\hfil}}}%
    {\mathclose{\setbox0=\hbox{$\textstyle        \lrgrdel$}\hbox to\wd0
                         {\hfil$\textstyle        \rangle  $\hfil}}}%
    {\mathclose{\setbox0=\hbox{$\scriptstyle      \lrgrdel$}\hbox to\wd0
                         {\hfil$\scriptstyle      \rangle  $\hfil}}}%
    {\mathclose{\setbox0=\hbox{$\scriptscriptstyle\lrgrdel$}\hbox to\wd0
                         {\hfil$\scriptscriptstyle\rangle  $\hfil}}}}%
\def\parldel {\mathchoice
    {\mathopen {\setbox0=\hbox{$\displaystyle     \lrgldel$}\hbox to\wd0
                         {\hfil$\displaystyle     [       $\hfil}}}%
    {\mathopen {\setbox0=\hbox{$\textstyle        \lrgldel$}\hbox to\wd0
                         {\hfil$\textstyle        [        $\hfil}}}%
    {\mathopen {\setbox0=\hbox{$\scriptstyle      \lrgldel$}\hbox to\wd0
                         {\hfil$\scriptstyle      [        $\hfil}}}%
    {\mathopen {\setbox0=\hbox{$\scriptscriptstyle\lrgldel$}\hbox to\wd0
                         {\hfil$\scriptscriptstyle[        $\hfil}}}}%
\def\parrdel {\mathchoice
    {\mathclose{\setbox0=\hbox{$\displaystyle     \lrgrdel$}\hbox to\wd0
                         {\hfil$\displaystyle     ]       $\hfil}}}%
    {\mathclose{\setbox0=\hbox{$\textstyle        \lrgrdel$}\hbox to\wd0
                         {\hfil$\textstyle        ]        $\hfil}}}%
    {\mathclose{\setbox0=\hbox{$\scriptstyle      \lrgrdel$}\hbox to\wd0
                         {\hfil$\scriptstyle      ]        $\hfil}}}%
    {\mathclose{\setbox0=\hbox{$\scriptscriptstyle\lrgrdel$}\hbox to\wd0
                         {\hfil$\scriptscriptstyle]        $\hfil}}}}%

\def\eightpoint{\small}                         
%%%%%%%%%%%%%%%%%%%%%%%%%
\def\pluldel {\mathchoice
   {\mathopen {\setbox0=\hbox{$\displaystyle     \lrgldel$}\hbox to\wd0
                        {\hfil$\displaystyle     [       $\hfil}%
                        \kern-\wd0\hbox to\wd0
                        {\hss$\vcenter{\hbox{\eightpoint$\scriptscriptstyle\bullet$}}$\hss}}}%
   {\mathopen {\setbox0=\hbox{$\textstyle        \lrgldel$}\hbox to\wd0
                        {\hfil$\textstyle        [       $\hfil}%
                        \kern-\wd0\hbox to\wd0
                        {\hss$\vcenter{\hbox{\eightpoint$\scriptscriptstyle\bullet$}}$\hss}}}%
   {\mathopen {\setbox0=\hbox{$\scriptstyle      \lrgldel$}\hbox to\wd0
                        {\hfil$\scriptstyle      [       $\hfil}%
                        \kern-\wd0\hbox to\wd0
                        {\hss$\raise.1ex\hbox{\eightpoint$\scriptscriptstyle\bullet$}$\hss}}}%
   {\mathopen {\setbox0=\hbox{$\scriptscriptstyle\lrgldel$}\hbox to\wd0
                        {\hfil$\scriptscriptstyle[       $\hfil}%
                        \kern-\wd0\hbox to\wd0
                        {\hss$\raise.03ex\hbox{\eightpoint$\scriptscriptstyle\bullet$}$\hss}}}}%
\def\plurdel {\mathchoice
   {\mathclose{\setbox0=\hbox{$\displaystyle     \lrgldel$}\hbox to\wd0
                        {\hfil$\displaystyle     ]       $\hfil}%
                        \kern-\wd0\hbox to\wd0
                        {\hss$\vcenter{\hbox{\eightpoint$\scriptscriptstyle\bullet$}}$\hss}}}%
   {\mathclose{\setbox0=\hbox{$\textstyle        \lrgldel$}\hbox to\wd0
                        {\hfil$\textstyle        ]       $\hfil}%
                        \kern-\wd0\hbox to\wd0
                        {\hss$\vcenter{\hbox{\eightpoint$\scriptscriptstyle\bullet$}}$\hss}}}%
   {\mathclose{\setbox0=\hbox{$\scriptstyle      \lrgldel$}\hbox to\wd0
                        {\hfil$\scriptstyle      ]       $\hfil}%
                        \kern-\wd0\hbox to\wd0
                        {\hss$\raise.1ex\hbox{\eightpoint$\scriptscriptstyle\bullet$}$\hss}}}%
   {\mathclose{\setbox0=\hbox{$\scriptscriptstyle\lrgldel$}\hbox to\wd0
                        {\hfil$\scriptscriptstyle]       $\hfil}%
                        \kern-\wd0\hbox to\wd0
                        {\hss$\raise.03ex\hbox{\eightpoint$\scriptscriptstyle\bullet$}$\hss}}}}%
\def\witldel {\mathchoice
   {\mathopen {\setbox0=\hbox{$\displaystyle     \lrgldel$}\hbox to\wd0
                        {\hfil$\displaystyle     (       $\hfil}%
                        \kern-\wd0\hbox to\wd0
                        {\hss$\vcenter{\hbox{\eightpoint$\scriptscriptstyle\bullet\mkern3.2mu$}}$\hss}}}%
   {\mathopen {\setbox0=\hbox{$\textstyle        \lrgldel$}\hbox to\wd0
                        {\hfil$\textstyle        (       $\hfil}%
                        \kern-\wd0\hbox to\wd0
                        {\hss$\vcenter{\hbox{\eightpoint$\scriptscriptstyle\bullet\mkern3.2mu$}}$\hss}}}%
   {\mathopen {\setbox0=\hbox{$\scriptstyle      \lrgldel$}\hbox to\wd0
                        {\hfil$\scriptstyle      (       $\hfil}%
                        \kern-\wd0\hbox to\wd0
                        {\hss$\raise.1ex\hbox{\eightpoint$\scriptscriptstyle\bullet\mkern3.2mu$}$\hss}}}%
   {\mathopen {\setbox0=\hbox{$\scriptscriptstyle\lrgldel$}\hbox to\wd0
                        {\hfil$\scriptscriptstyle(       $\hfil}%
                        \kern-\wd0\hbox to\wd0
                        {\hss$\raise.03ex\hbox{\eightpoint$\scriptscriptstyle\bullet\mkern3.2mu$}$\hss}}}}%
\def\witrdel {\mathchoice
   {\mathclose{\setbox0=\hbox{$\displaystyle     \lrgldel$}\hbox to\wd0
                        {\hfil$\displaystyle     )       $\hfil}%
                        \kern-\wd0\hbox to\wd0
                        {\hss$\vcenter{\hbox{\eightpoint$\scriptscriptstyle\mkern3.2mu\bullet$}}$\hss}}}%
   {\mathclose{\setbox0=\hbox{$\textstyle        \lrgldel$}\hbox to\wd0
                        {\hfil$\textstyle        )       $\hfil}%
                        \kern-\wd0\hbox to\wd0
                        {\hss$\vcenter{\hbox{\eightpoint$\scriptscriptstyle\mkern3.2mu\bullet$}}$\hss}}}%
   {\mathclose{\setbox0=\hbox{$\scriptstyle      \lrgldel$}\hbox to\wd0
                        {\hfil$\scriptstyle      )       $\hfil}%
                        \kern-\wd0\hbox to\wd0
                        {\hss$\raise.1ex\hbox{\eightpoint$\scriptscriptstyle\mkern3.2mu\bullet$}$\hss}}}%
   {\mathclose{\setbox0=\hbox{$\scriptscriptstyle\lrgldel$}\hbox to\wd0
                        {\hfil$\scriptscriptstyle)       $\hfil}%
                        \kern-\wd0\hbox to\wd0
                        {\hss$\raise.03ex\hbox{\eightpoint$\scriptscriptstyle\mkern3.2mu\bullet$}$\hss}}}}%
%%%%%%%%%%%%%%%%%%%%%%%%%%%%

\newbox\ldelbox
\setbox\ldelbox=\hbox{$\lrgldel$}
\def\ldelskip{\hskip\wd\ldelbox}
\newbox\rdelbox
\setbox\rdelbox=\hbox{$\lrgrdel$}
\def\rdelskip{\hskip\wd\rdelbox}

\def\wits #1{\witldel #1\witrdel}%                         
\def\plus #1{\pluldel #1\plurdel}%
                         
\def\aprs #1{\aprldel #1\aprrdel}%
\def\pars #1{\parldel #1\parrdel}%
\def\seqs #1{\seqldel #1\seqrdel}%
\def\cons #1{\conldel #1\conrdel}%
\def\copt #1{#1}%



%% linear logic %%%%%%%%%%%%%%%%%%%%%%%%%%%%

\let\turnstile=\vdash
\def\lone{1}
\def\lbot{\bot}
\def\lzero{0}
\def\ltop{\top}
\def\ltens{\mathop{\varotimes}}
\def\lcut{\mathop\varobar}
\def\lpar{\mathop\bindnasrepma}
\def\lplus{\mathop\varoplus}
\def\lwith{\mathop\binampersand}
\def\lwn{\mathord{?}}
\def\loc{\mathord{!}}
\def\lneg{^\bot}
\def\lnegneg{^{\bot\bot}}
\def\limp{\multimap}
\def\lseq{\mathop\vartriangleleft}
\def\lcoseq{\mathop\vartriangleright}

\def\lweakr{\mathord{\lwn}\mathsf{w}}
\def\lconr{\mathord{\lwn}\mathsf{c}}
\def\lderr{\mathord{\lwn}\mathsf{d}}

%%%%%%%%%%%%%%%%%%%%%%%%%%%%%%%%%%%%%%%%%%%%%%%%%%%%%%%%%%%
% sonstiges
 
%% \newcommand{\ie}{i.e.,\ }
%% \newcommand{\eg}{e.g.,\ }
%% \newcommand{\cf}{cf.\ }
%% \newcommand{\etc}{etc.\ }
%% \newcommand{\etal}{et al.\ }
%% \newcommand{\Wolg}{Without loss of generality}
%% \newcommand{\wolg}{without loss of generality}
%% \newcommand{\bwoc}{by way of contradiction}
%% \newcommand{\Bwoc}{By way of contradiction}

\def\quadfs {\rlap{\rm\quad.}}%
\def\quadcm {\rlap{\rm\quad,}}%
\def\quadsc {\rlap{\rm\quad;}}%
\def\quadcl {\rlap{\rm\quad:}}%
\def\quadqm {\rlap{\rm\quad?}}%
\def\quand {\quad\mbox{and}\quad}%
\def\qquand {\qquad\mbox{and}\qquad}%
\def\qqquand {\quad\qquad\mbox{and}\qquad\quad}%
\def\qqqquand {\qquad\qquad\mbox{and}\qquad\qquad}%
\def\quor {\quad\mbox{or}\quad}%
\def\qquor {\qquad\mbox{or}\qquad}%
\def\qqquor {\quad\qquad\mbox{or}\qquad\quad}%
\def\qqqquor {\qquad\qquad\mbox{or}\qquad\qquad}%

\def\quby {\quad\mbox{by}\quad}%
\def\qquby {\qquad\mbox{by}\qquad}%
\def\qqquby {\quad\qquad\mbox{by}\qquad\quad}%
\def\qqqquby {\qquad\qquad\mbox{by}\qquad\qquad}%

\def\quato {\quad\to\quad}%
\def\qquato {\qquad\to\qquad}%
\def\qqquato {\quad\qquad\to\qquad\quad}%
\def\qqqquato {\qquad\qquad\to\qquad\qquad}%

\def\dqualto {\hskip.5em\leadsto\hskip.5em}%
\def\qualto {\quad\leadsto\quad}%
\def\qqualto {\qquad\leadsto\qquad}%
\def\qqqualto {\quad\qquad\leadsto\qquad\quad}%
\def\qqqqualto {\qquad\qquad\leadsto\qquad\qquad}%

\def\quand {\quad\mbox{and}\quad}%
\def\qquand {\qquad\mbox{and}\qquad}%
\def\qqquand {\quad\qquad\mbox{and}\qquad\quad}%
\def\qqqquand {\qquad\qquad\mbox{and}\qquad\qquad}%

\def\nquand {\quad\mbox{\normalsize and}\quad}%
\def\nqquand {\qquad\mbox{\normalsize and}\qquad}%
\def\nqqquand {\quad\qquad\mbox{\normalsize and}\qquad\quad}%
\def\nqqqquand {\qquad\qquad\mbox{\normalsize and}\qquad\qquad}%

\def\quor {\quad\mbox{or}\quad}%
\def\qquor {\qquad\mbox{or}\qquad}%
\def\qqquor {\quad\qquad\mbox{or}\qquad\quad}%
\def\qqqquor {\qquad\qquad\mbox{or}\qquad\qquad}%

\def\quwith {\quad\mbox{with}\quad}%
\def\qquwith {\qquad\mbox{with}\qquad}%
\def\qqquwith {\quad\qquad\mbox{with}\qquad\quad}%
\def\qqqquwith {\qquad\qquad\mbox{with}\qquad\qquad}%

\def\quiff {\quad\mbox{iff}\quad}%
\def\qquiff {\qquad\mbox{iff}\qquad}%
\def\qqquiff {\quad\qquad\mbox{iff}\qquad\quad}%
\def\qqqquiff {\qquad\qquad\mbox{iff}\qquad\qquad}%

\def\quifof {\quad\mbox{if and only if}\quad}%
\def\qquifof {\qquad\mbox{if and only if}\qquad}%
\def\qqquifof {\quad\qquad\mbox{if and only if}\qquad\quad}%
\def\qqqquifof {\qquad\qquad\mbox{if and only if}\qquad\qquad}%

\def\queq {\quad=\quad}%
\def\qqueq {\qquad=\qquad}%

\def\qquadfs {\rlap{\rm\qquad.}}%
\def\qquadcm {\rlap{\rm\qquad,}}%

\def\yields{\quad\mbox{yields}\quad}

\def\mydisplaybreak{\hspace*{\fill} \cr \hspace*{\fill}}

\def\vcvdots{\vbox{\baselineskip=4pt \lineskiplimit=0pt
    \hbox{.}\hbox{.}\hbox{.}}}


%  \clap and \proofadjust ------------------------

\def\nlap#1{#1}
\def\clap#1{\hbox to 0pt{\hss#1\hss}}
\def\sqlap#1{\hbox to .5em{\hss#1\hss}}
\def\qlap#1{\hbox to 1em{\hss#1\hss}}
\def\qqlap#1{\hbox to 2em{\hss#1\hss}}
\def\qqqlap#1{\hbox to 3em{\hss#1\hss}}
\def\qqqqlap#1{\hbox to 4em{\hss#1\hss}}
\def\qqqqqlap#1{\hbox to 5em{\hss#1\hss}}
\def\qqqqqqlap#1{\hbox to 6em{\hss#1\hss}}
\def\qqqqqqqlap#1{\hbox to 7em{\hss#1\hss}}
\def\qqqqqqqqlap#1{\hbox to 8em{\hss#1\hss}}
\def\qqqqqqqqqlap#1{\hbox to 9em{\hss#1\hss}}
\newcommand{\wlap}[2][10ex]{\hbox to#1{\hss#2\hss}}
\def\nlapm#1{#1}
\def\clapm#1{\clap{$#1$}}
\def\sqlapm#1{\sqlap{$#1$}}
\def\qlapm#1{\qlap{$#1$}}
\def\qqlapm#1{\qqlap{$#1$}}
\def\qqqlapm#1{\qqqlap{$#1$}}
\def\qqqqlapm#1{\qqqqlap{$#1$}}
\def\qqqqqlapm#1{\qqqqqlap{$#1$}}
\def\qqqqqqlapm#1{\qqqqqqlap{$#1$}}
\def\qqqqqqqlapm#1{\qqqqqqqlap{$#1$}}
\def\qqqqqqqqlapm#1{\qqqqqqqqlap{$#1$}}
\def\qqqqqqqqqlapm#1{\qqqqqqqqqlap{$#1$}}
\newcommand{\wlapm}[2][10ex]{\hbox to#1{\hss$#2$\hss}}
\def\rlapm#1{\hbox to 0pt{$#1$\hss}}
\def\llapm#1{\hbox to 0pt{\hss$#1$}}

\def\squad{\hskip.5em}
\def\qqquad{\quad\qquad}
\def\qqqquad{\qquad\qquad}
\def\qqqqquad{\qqquad\qquad}
\def\qqqqqquad{\qqquad\qqquad}

\newcommand{\vclap}[2][0pt]{\hbox to #1{\hss#2\hss}}
\newcommand{\vclapm}[2][0pt]{\hbox to #1{\hss$#2$\hss}}



\def\proofadjust{\vadjust{\nobreak\vskip-2.7ex\nobreak}}
\def\semiproofadjust{\vadjust{\nobreak\vskip-1.3ex\nobreak}}
\def\doubleproofadjust{\vadjust{\nobreak\vskip-4ex\nobreak}}
\def\tripleproofadjust{\vadjust{\nobreak\vskip-5.3ex\nobreak}}


\def\interdisplayskip{.5ex}
\newskip\mydisplaywidth
\newcommand{\twolinedisplay}[3][10pt]{%
  \mydisplaywidth=\displaywidth
  \advance\mydisplaywidth-#1
  \begin{array}{c}
    \clap{\hbox to\mydisplaywidth{$\displaystyle#2$\hss}}\\[\interdisplayskip]
    \clap{\hbox to\mydisplaywidth{\hss$\displaystyle#3$}}
  \end{array}
}

\def\interarrayskip{.5ex}

\def\nosmash#1{#1}
%%%%%%%%%%%%%%%%%%%%%%%%%%%%%%%%%%%%%%%%%%%%%%%%%%

\endinput


\DeclareMathSymbol{:}{\mathbin}{operators}{"3A}


\usepackage{graphicx}
\usepackage{amstext}
\usepackage{wasysym}


\newcommand*{\Cdot}[1][2.5]{%
  \mathpalette{\CdotAux{#1}}\cdot%
}
\newdimen\CdotAxis
\newcommand*{\CdotAux}[3]{%
  {%
    \settoheight\CdotAxis{$#2\vcenter{}$}%
    \sbox0{%
      \raisebox\CdotAxis{%
        \scalebox{#1}{%
          \raisebox{-\CdotAxis}{%
            $\mathsurround=0pt #2#3$%
          }%
        }%
      }%
    }%
    % Remove depth that arises from scaling.
    \dp0=0pt %
    % Decrease scaled height.
    \sbox2{$#2\bullet$}%
    \ifdim\ht2<\ht0 %
      \ht0=\ht2 %
    \fi
    % Use the same width as the original \cdot.
    \sbox2{$\mathsurround=0pt #2#3$}%
    \hbox to \wd2{\hss\usebox{0}\hss}%
  }%
}


\newcommand{\Vars}{\mathit{Vars}}

\newcommand{\md}{\mathsf{md}}
\newcommand{\dmd}{\mathsf{dmd}}

% \renewcommand\merge[2]{%
%   \ooalign{\hfil$\vcenter{\hbox{$#1$}}$\hfil\cr
%     \hfil$\vcenter{\hbox{$#2$}}$\hfil}}

%\newcommand{\dyndiamond}{\raisebox{.4mm}{\merge\lozenge\Cdot}}
%\newcommand{\dynbox}{\raisebox{.4mm}{\merge\square\Cdot}}


\newcommand{\dyndiamond}[1]{{\fontsize{7}{7}\selectfont\tup{\raisebox{.6pt}{\textsf{#1}}}}}
\newcommand{\dynbox}[1]{{\fontsize{7}{7}\selectfont[\raisebox{.6pt}{\textsf{#1}}]}}

\newcommand\dswap{\dyndiamond{lsw}}
\newcommand\bswap{\dynbox{lsw}}
\newcommand\msw{\dyndiamond{lsw}}
\newcommand\mgsw{\dyndiamond{gsw}}

\newcommand\sab{\dyndiamond{gsb}}
\newcommand\sabbox{\dynbox{gsb}}
\newcommand\msb{\dyndiamond{lsb}}
\newcommand\mgsb{\dyndiamond{gsb}}

\newcommand{\lsab}{\dyndiamond{lsb}}
\newcommand{\lsabbox}{\dynbox{lsb}}

\newcommand{\bridge}{\dyndiamond{lbr}}
\newcommand{\bridgebox}{\dynbox{lbr}}
\newcommand\mbr{\dyndiamond{lbr}}
\newcommand\mgbr{\dyndiamond{gbr}}


\newcommand\dgswap{\dyndiamond{gsw}}
\newcommand\bgswap{\dynbox{gsw}}

\newcommand{\gbridge}{\dyndiamond{gbr}}
\newcommand{\gbridgebox}{\dynbox{gbr}}

\newcommand{\trbr}{\Circle^+}
\newcommand{\kn}{\textcircled{\small{k}}}
\newcommand{\re}{\textcircled{r}}

\newcommand{\qbfmodels}{\models_{\mathrm{qbf}}}

\newcommand{\iml}{\ensuremath{\mathsf{IML}}\xspace}
\newcommand{\bml}{\ensuremath{\mathcal{ML}}\xspace}
\newcommand{\bhl}{\ensuremath{\mathcal{HL}}\xspace}
\newcommand{\swl}{{\bml(\dswap)}}
\newcommand{\gsl}{{\bml(\sab)}}
\newcommand{\lsl}{{\bml(\lsab)}}
\newcommand{\brl}{{\bml(\bridge)}}
\newcommand{\gswl}{{\bml(\dgswap)}}
\newcommand{\gbrl}{{\bml(\gbridge)}}

% ---Memory Logics
\newcommand{\mlo}{{\cal ML}(\re,\kn)}

\newcommand{\trbrl}{{\bml(\trbr)}}

\newcommand{\fol}{{\mathcal{FOL}}}
%\newcommand{\sol}{{\mathcal{SOL}}}
\newcommand{\pow}{{\mathcal{P}}}

\newcommand{\down}{{\downarrow}}
\newcommand{\Hl}{\mathcal{H}}
\newcommand{\hlsd}{\Hl({:},\down)}
\newcommand{\hlsde}{\Hl({:},\down,\E)}

\newcommand{\A}{\textsf{A}}
\newcommand{\E}{\textsf{E}}
\newcommand{\PREF}{\mathsf{PREF}}
\newcommand{\DYN}{{\rm \sf DYN}\xspace}
\newcommand{\PROP}{{\rm \sf PROP}\xspace}
\newcommand{\VAR}{{\rm \sf VAR}\xspace}
\newcommand{\NOM}{{\rm \sf NOM}\xspace}
\newcommand{\FORM}{{\rm \sf FORM}\xspace}
\newcommand{\REL}{{\rm \sf REL}\xspace}
\newcommand{\ra}{\rightarrow}
\newcommand{\lra}{\leftrightarrow}
\newcommand{\Tr}{{\sf MTr}}
%\newcommand{\bisim}{\leftrightarroweq}
\renewcommand{\iff}{\mbox{iff}}
\newcommand{\eqrel}{\leftrightsquigarrow}
\renewcommand{\L}{{\cal L}}


%%%% Aca empiezan las macros viejas

\newcommand{\tup}[1]{\langle #1 \rangle}

\newcommand{\ttup}[1]{\langle\!\langle #1 \rangle\!\rangle}

\newcommand{\model}{\ensuremath{\mathcal{M}}\xspace}
\newcommand{\nodel}{\ensuremath{\mathcal{N}}\xspace}
\newcommand{\univ}{\ensuremath{\mathcal{U}}\xspace}
\newcommand{\uset}{U}
\newcommand{\img}{\textrm{im\ }}
\newcommand{\bisim}{\mathrel{\,\raisebox{.3ex}{$\underline{\makebox[.7em]{$\leftrightarrow$}}$}\,}}
\newcommand{\tbisim}{\bisim^T}

\newcommand{\cset}[1]{\{#1\}}
\newcommand{\csetsc}[2]{\{ #1 \mid #2\}}

\newcommand{\gM}{\mathcal{M}}
\newcommand{\gC}{\mathcal{C}}
\newcommand{\gL}{\mathcal{L}}


%%%%%%%%% de facu
% \theoremstyle{definition}
% \newtheorem{theorem}{Theorem}[section]
% \newtheorem{lemma}[theorem]{Lemma}
% \newtheorem{proposition}[theorem]{Proposition}
% \newcounter{coro}
% \newtheorem{corollary}[theorem]{Corollary}
% \newtheorem{conjecture}[theorem]{Conjecture}
% % \newtheorem{proof}{Proof}

% \newtheorem{obs}[theorem]{Observation}
% \newtheorem{definition}[theorem]{Definition}
% \newtheorem{construction}[theorem]{Construction}
% \newtheorem{nota}[theorem]{Notation}
% \newtheorem{example}[theorem]{Example}

%\newtheorem*{thm*}{Theorem} % sin numeracion
%\newtheorem*{defn*}{Definition}


\newcommand{\theq}[1]{\equiv_{#1}}

\newcommand{\elementary}{\ensuremath{EC_{\Delta}}\xspace}
\newcommand{\belementary}{\ensuremath{EC}\xspace}
\newcommand{\setcomp}[2]{\{{#1} : {#2}\}}

\newcommand{\diam}[1]{\langle #1 \rangle}
\newcommand{\ddiam}[1]{\ttup{#1}}
\newcommand{\ST}{\mathsf{STr}}

\newcommand{\HT}{\mathsf{HTr}}

%%%%%%%%%%%%%%%%%%%%%%%%%%%%%%%%%%%%%%
%---------Complexity Classes---------%
%%%%%%%%%%%%%%%%%%%%%%%%%%%%%%%%%%%%%%
\newcommand{\ntime}{\textsc{NTime}\xspace}
\newcommand{\ptime}{\textsc{P}\xspace}
\newcommand{\nptime}{\textsc{NP}\xspace}
\newcommand{\nlogspace}{\textsc{NLogSpace}\xspace}
\newcommand{\pspace}{\textsc{PSpace}\xspace}
\newcommand{\npspace}{\textsc{NPSpace}\xspace}
\newcommand{\exptime}{\textsc{ExpTime}\xspace}
\newcommand{\nexptime}{\textsc{NExpTime}\xspace}
\newcommand{\expspace}{\textsc{ExpSpace}\xspace}
\newcommand{\order}[1]{{\cal O} (#1)}
\newcommand{\C}{{\cal C}}

\newcommand{\wrv}{\tup{W,R,V}}
\newcommand{\wrvp}{\tup{W',R',V'}}

\newcommand{\EF}{\mathit{EF}}
\newcommand{\sbt}{\,\begin{picture}(-1,1)(-1,-3)\circle*{3}\end{picture}\ }


\newcommand{\sml}{\bml(\sdia)}
\newcommand{\shl}{\bhl(\sdia)}

\usepackage{enumitem}
\usepackage{graphicx}
\usepackage{multicol}
\usepackage{xargs}
\usepackage{tikz-cd}
\usepackage{amssymb}
\usepackage{vwcol}

\usepackage{tikz}
\usetikzlibrary{arrows,calc,patterns,positioning,shapes}
\usetikzlibrary{decorations.pathmorphing}
\tikzset{
    modal/.style={>=stealth',shorten >=1pt,shorten <=1pt,auto,
    node distance=1.5cm,semithick},
    world/.style={circle,draw,minimum size=.8cm,fill=gray!15},
    point/.style={circle,draw,fill=black,inner sep=0.5mm},
    reflexive/.style={->,in=120,out=60,loop,looseness=#1},
    reflexive/.default={5},
    reflexive point/.style={->,in=135,out=45,loop,looseness=#1},
    reflexive point/.default={25},
}

\newcommand{\valentin}[1]{{\color{violet}[[V: #1]]}}

%%% Math Euler (fonts) %%%

% \DeclareMathAlphabet\EuRoman{U}{eur}{m}{n}
% \SetMathAlphabet\EuRoman{bold}{U}{eur}{b}{n}
% \newcommand{\matheu}{\EuRoman}

%%% Kripke Models

\newcommand{\props}{\mathsf{Prop}}

\newcommand{\W}{\mathrm{Q}}
\newcommand{\R}{\mathbin{\mathrm{R}}}
\newcommand{\V}{\mathrm{V}}

\newcommand{\Q}{\mathrm{Q}}

\newcommand{\Wof}[1]{\mathrm{Q}_{#1}}
\newcommand{\Qof}[1]{\mathrm{Q}_{#1}}
\newcommand{\Rof}[1]{\mathbin{\mathrm{R}_{#1}}}
\newcommand{\iRof}[1]{\mathbin{\mathrm{R}^{-1}_{#1}}}
\newcommand{\Vof}[1]{\mathrm{V}_{#1}}

\newcommand{\kmodel}{\mathfrak{K}}
\newcommand{\ikmodel}{\mathfrak{K}^{-1}}

\newcommand{\Wkm}[1][]{\Wof{\kmodel_{#1}}}
\newcommand{\Qkm}[1][]{\Qof{\kmodel_{#1}}}
\newcommand{\Rkm}[1][]{\Rof{\kmodel_{#1}}}
\newcommand{\iRkm}[1][]{\iRof{\kmodel_{#1}}}
\newcommand{\Vkm}[1][]{\Vof{\kmodel_{#1}}}

%%%% Frames

% \newcommand{\kframe}{\matheu{F}}

% \newcommand{\Wfm}[1][]{\Wof{\kframe_{#1}}}
% \newcommand{\Rfm}[1][]{\Rof{\kframe_{#1}}}
% \newcommand{\Vfm}[1][]{\Vof{\kframe_{#1}}}

% \newcommand{\fmofkm}{\kframe(\kmodel)}
% \newcommandx{\kmoffmv}[2][1={\kframe},2={\V}]{\kmodel(\npair {#1} {#2})}

%%%% Filtrations

\newcommand{\fset}{\Sigma}
\newcommandx{\fksym}[2][1=\fset,2=\kmodel]{^{#1}_{#2}}
\newcommandx{\feqr}[2][1=\kmodel,2=\fset]{\mathrel{{\equiv}{\fksym[#1][#2]}}}
\newcommandx{\feqc}[3][2=\kmodel,3=\fset]{[#1]\fksym[{#3}][{#2}]}

\newcommand{\fmodel}[1][]{\kmodel_{#1}^{\fset}}

\newcommand{\Wfm}[1][]{\Wof{\kmodel_{#1}}^{\fset}}
\newcommand{\Qfm}[1][]{\Qof{\kmodel_{#1}}^{\fset}}
\newcommand{\Rfm}[1][]{\R_{\kmodel_{#1}}^{\fset}}
\newcommand{\Vfm}[1][]{\Vof{\kmodel_{#1}}^{\fset}}


%%%% Bi-relational Models

\usepackage{xargs}
\newcommandx{\bmodel}[1][1={B}]{\mathfrak{#1}}

\newcommand{\bmsymb}{\mathfrak{B}}

\newcommand{\Wbm}[1][]{\Wof{\bmsymb_{#1}}}
\newcommand{\Qbm}[1][]{\Wof{\bmsymb_{#1}}}
\newcommand{\Rbm}[1][]{\Rof{\bmsymb_{#1}}}
\newcommand{\Vbm}[1][]{\Vof{\bmsymb_{#1}}}

\newcommand{\Fsymb}{\preccurlyeq}
\newcommand{\Fbm}{{\Fsymb_\bmsymb}}

\newcommand{\kmofbm}{\kmodel(\bmodel)}
\newcommand{\bmofhm}{\bmodel(\hmodel)}
\newcommand{\bmofwm}{\bmodel(\wmodel)}

\newcommand{\bsat}{\mathrm{b}}

%%% Hierarchical Kripke Models

\newcommand{\Kof}[1]{\mathcal{K}_{#1}}
\newcommand{\Kmodels}{\Kof{}}

\newcommand{\Fof}[1]{\Fsymb_{#1}}

\newcommand{\imodel}{\mathfrak{I}}
\newcommand{\hmodel}{\mathfrak{H}}

\newcommand{\Hleqsym}{\preccurlyeq}

\newcommand{\Hleqof}[1][]{\Hleqsym_{#1}}
\newcommand{\Hleq}{\Hleqof[\hmodel]}
\newcommand{\Fhm}[1][]{\Fof{\hmodel_{#1}}}

\newcommand{\Khm}[1][]{\Kof{\hmodel_{#1}}}

%%% Weaved Kripke Models

\newcommand{\wmodel}{\mathfrak{W}}
\newcommand{\tmodel}{\mathfrak{T}}

\newcommand{\fsim}{\rightsquigarrow}
\newcommand{\BSim}{\leftrightarroweq}

\newcommand{\Sim}{\mathbin{\mathsf{S}}}
\newcommand{\Tim}{\mathbin{\mathsf{T}}}

% \usepackage{mathrsfs}
% \newcommand{\G}{\mathcal{S}}
\newcommand{\Sof}[1]{\mathcal{S}_{#1}}
\newcommand{\Sims}{\Sof{}}

\newcommand{\Kwm}[1][]{\Kof{\wmodel_{#1}}}
\newcommand{\Swm}[1][]{\Sof{\wmodel_{#1}}}
\newcommand{\Fwm}[1][]{{\Fsymb_{\wmodel_{#1}}}}
\newcommand{\Qwm}[1][]{\Qof{\wmodel_{#1}}}

\newcommand{\Ktm}[1][]{\Kof{\tmodel_{#1}}}
\newcommand{\Stm}[1][]{\Sof{\tmodel_{#1}}}
\newcommand{\Ftm}[1][]{{\Fsymb_{\tmodel_{#1}}}}
\newcommand{\Qtm}[1][]{\Qof{\tmodel_{#1}}}

\newcommand{\fwmodel}{\wmodel^{\fset}}
\newcommand{\Kfwm}{\Kmodels^{\fset}}
\newcommand{\Sfwm}{\Sims^{\fset}}
\newcommand{\fSim}{\Sim^{\fset}}

\newcommand{\rwmodel}{\restr{\wmodel}{\kmodel}}
\newcommand{\Krwm}{\restr{\Kmodels}{\kmodel}}
\newcommand{\Srwm}{\restr{\Sims}{\kmodel}}

\newcommand{\F}{\mathcal{F}}
\newcommand{\uwmodel}{\restr{\wmodel}{\kmodel^{*}}}
\newcommand{\Kuwm}{\restr{\Kmodels}{\kmodel^{*}}}
\newcommand{\Suwm}{\restr{\Sims}{\kmodel^{*}}}

\newcommand{\pwmodel}{\restr{\wmodel}{\kmodel^{n}}}
\newcommand{\Kpwm}{\restr{\Kmodels}{\kmodel^{n}}}
\newcommand{\Spwm}{\restr{\Sims}{\kmodel^{n}}}


\newcommand{\branches}{\mathcal{C}}
\newcommand{\bwmodel}{\mathfrak{C}}
\newcommand{\Kbwm}[1][]{\Kof{\bwmodel_{#1}}}
\newcommand{\Sbwm}[1][]{\Sof{\bwmodel_{#1}}}

\newcommand{\twmodel}{\mathfrak{T}}
\newcommand{\Ktwm}{\Kof{\twmodel}}
\newcommand{\Stwm}{\Sof{\twmodel}}

\newcommand{\iwmodel}{\wmodel^{\cong}}
\newcommand{\Kiwm}{\Kmodels^{\cong}}
\newcommand{\Siwm}{\Sims^{\cong}}
\newcommand{\iSim}{\Sim^{\cong}}
\newcommand{\IsoKiwm}{\mathcal{I}^{\cong}}

\newcommand{\cwmodel}{\wmodel^{\Downarrow}}
\newcommand{\Kcwm}{\Kmodels^{\Downarrow}}
\newcommand{\Scwm}{\Sims^{\Downarrow}}

%%% Weaved Frames

% \newcommand{\wframe}{\matheu{W}}

% \newcommand{\Kwf}[1][]{\mathbb{F}_{\wframe_{#1}}}
% \newcommand{\Swf}[1][]{\Sof{\wframe_{#1}}}
% \newcommand{\Vwf}[1][]{\V_{#1}}

% \newcommand{\Qwf}[1][]{\Qof{\wframe_{#1}}}
% \newcommand{\Fwf}{{\Fsymb_\wframe}}
% \newcommand{\nFwf}{{{\not\Fsymb}_\wframe}}

%%%% From ... 2 Weaved ...

\newcommand{\QQs}[1][]{\mathcal{Q}_{#1}}
\newcommand{\Ind}{I}
\newcommand{\bigq}[1][]{\bm{Q_{#1}}}

\newcommand{\wmofhm}{\wmodel(\hmodel)}
\newcommand{\wmofbm}[1]{\wmodel_{#1}(\bmodel)}
\newcommand{\wmofwfv}{\wmodel(\wframe, \Vwf)}
\newcommand{\wfofwm}{\wmodel(\Vwf)}

%%% MISC

\usepackage{mathtools}
\DeclarePairedDelimiterX\setof[2]{\{}{\}}{\,#1 \;\delimsize\vert\; #2\,}

\DeclareMathOperator{\dom}{dom}
\DeclareMathOperator{\cod}{cod}

\newcommand{\last}[1]{#1_{{-1}}}

\newcommand{\imap}{\iota}
\newcommand{\isym}{\hookrightarrow}  % inclusion function
\newcommand{\bsym}{%
    % \hookrightarrow\mathrel{\mspace{-15mu}}\rightarrow
    % \mathbin{\tikz{\draw[>->>] (0,0) -- (0.4,0);}}
    \cong
} % bijection function

\DeclareMathOperator{\ideg}{deg}
\DeclareMathOperator{\size}{size}

\newcommand{\restr}[2]{#1\vert_{#2}}

% \DeclareMathOperator{\id}{id}

\newcommand{\bb}[1]{\mathbf{#1}}

%%% REFERENCING

\usepackage{hyperref}
\usepackage{cleveref}

\Crefname{definition}{Def.}{Defs.}
\Crefname{figure}{Fig.}{Figs.}
\Crefname{proposition}{Prop.}{Props.}
\Crefname{equation}{Eq.}{Eqs.}
\Crefname{property}{Prop.}{Props.}
\Crefname{theorem}{Thm.}{Thms.}
\Crefname{thm}{Thm.}{Thms.}
\Crefname{example}{Ex.}{Exs.}
\Crefname{section}{Sec.}{Secs.}
\Crefname{fact}{Fact}{Facts}
\Crefname{appendix}{App.}{App.}
\Crefname{table}{Tab.}{Tabs.}

%%% INLINE EQUATIONS %%%
\newcounter{ieqnnum}
\makeatletter
\newcommand*{\ieqn}[2][]{%
    \begingroup
        % Put \refstepcounter at the beginning, because
        % package `hyperref' sets the anchor here.
        \refstepcounter{ieqnnum}%
        (\roman{ieqnnum})\xspace%
        \ifx\\#1\\%
            % \label[ieqn]{\theieqnnum}
        \else
            \label[ieqn]{#1}
        \fi
        % prevent line breaks inside equation
        \relpenalty=10000 %
        \binoppenalty=10000 %
        \ensuremath{%
            % \displaystyle % larger fractions, ...
            #2%
        }%
    \endgroup
}
\makeatother

\newcounter{rieqnnum}%
\newcommand{\ieqnref}[1]{%
    \setcounter{rieqnnum}{#1}%
    \setcounter{rieqnnum}{%
        \numexpr\value{ieqnnum}+\value{rieqnnum}%
    }%
    (\roman{rieqnnum})%
}

\colorlet{ctxcolor}{gray!70!white}
\colorlet{notgreen}{green!70!black}

\newcommand{\lat}[1]{{\color{NavyBlue}@_{#1}}}
\newcommand{\latr}[1]{@_{#1}}
\newcommand{\lbox}{\Box}
\newcommand{\ldia}{\Diamond}
\newcommand{\sbbox}{[\mathsf{s}]}
\newcommand{\sdia}{\langle\mathsf{s}\rangle}
\newcommand{\sloc}[1]{{\color{BlueViolet}(\mkern-6mu\{\mkern0mu #1 \mkern0mu\}\mkern-6mu)}} %%{\bigcirc^{#1}}
%\newcommand{\land}{\wedge}
%\newcommand{\lor}{\vee}
\renewcommand{\cimp}{\mathbin{\to}}

%%%% Connectives
\newcommand*{\NOT}{\neg}
\newcommand*{\AND}{\mathbin{\wedge}}
\newcommand*{\TOP}{\mathord{\top}}
\newcommand*{\OR}{\mathbin{\vee}}
\newcommand*{\BOT}{\mathord{\bot}}
\newcommand*{\IMP}{\mathbin{\supset}}%%\scalebox{.9}{\raise.2ex\hbox{$\supset$}}}}
\newcommand*{\BOX}{\mathord{\Box}}
\newcommand*{\DIA}{\mathord{\Diamond}}


\newcommand{\lats}[2]{\lat{#1}\sloc{#2}}
\newcommand{\latd}[2]{\lat{#1}\ldia{#2}}
\renewcommand{\cneg}[1]{\neg#1}
%\newcommand{\ssqn}[2]{#1\vdash#2}
\newcommand{\ssqnctx}[4]{\ssqn{{\color{ctxcolor} #1} #2}{#3 {\color{ctxcolor} #4}}}
\newcommand{\ssqnGD}[2]{\ssqnctx{\Gamma}{#1}{#2}{\Delta}}

\newcommand{\LR}[1]{{#1}_\mathsf{L}}
\newcommand{\RR}[1]{{#1}_\mathsf{R}}
  
\newcommand{\edge}[2]{\overline{#1#2}}

\newcommand{\proviso}[1]{\hbox{\small #1}}
\newcommand{\axr}{\mathsf{ax}}
\newcommand{\subar}{\mathsf{sub_1}}
\newcommand{\subbr}{\mathsf{sub_2}}

\newcommand{\lutz}[1]{{\color{notgreen}[[L: #1]]}} 
\newcommand{\carlos}[1]{{\color{red}[[C: #1]]}}
\newcommand{\todo}[1]{{\color{red}[[ToDo: #1]]}} 


\newcommand{\Layerssym}{\mathfrak{U}} %%\mathcal{L}}
\newcommand{\Layers}[1][]{\Layerssym_{#1}}
\newcommand{\layersym}{\mathfrak{u}} %%{L}
\newcommand{\layerasym}{\mathfrak{v}} %%{L}
\newcommand{\layer}{\layersym}
\newcommand{\layera}{\layerasym}
\newcommand{\lay}[1]{\layersym_{#1}}
\newcommand{\laya}[1]{\layerasym_{#1}}
\newcommand{\Worldssym}{W}
\newcommand{\Worlds}[1][\layer]{\Worldssym_{#1}}
\newcommand{\worldsym}{w}
\newcommand{\worldasym}{z}
\newcommand{\world}{\worldsym}
\newcommand{\worlda}{\worldasym}
\newcommand{\wor}[1]{\worldsym_{#1}}
\newcommand{\Relsym}{R}
\newcommand{\Rel}[1][\layer]{\Relsym_{#1}}
\newcommand{\Funsym}{F}
\newcommand{\Fun}[2]{\Funsym_{#1 \mapsto #2}}
\newcommand{\Funa}{\Fun{\layer}{\layera}}
\newcommand{\Valsym}{V}
\newcommand{\Val}[1][\layer]{\Valsym_{#1}}


\newcommand{\relcomp}{\mathord{;}}
%% \newcommand{\modelsym}{\mathcal{M}}
%% \newcommand{\model}{\modelsym}
\newcommand{\Framesym}{\mathcal{F}}
\newcommand{\Frame}{\Framesym}

\newcommand{\forcesym}{\Vdash}
\newcommand{\nforcesym}{\not\Vdash}

\newcommand{\forces}[2]{{#1}\forcesym{#2}}
\newcommand{\nforces}[2]{{#1}\nforcesym{#2}}

\newcommand{\xforces}[3]{\forces{#1,#2}{#3}}
\newcommand{\nxforces}[3]{\nforces{#1,#2}{#3}}

\newcommand{\bforces}[2]{\xforces{\bmodel}{#1}{#2}}
\newcommand{\nbforces}[2]{\nxforces{\bmodel}{#1}{#2}}

\newcommand{\hforces}[2]{\xforces{\hmodel}{#1}{#2}}
\newcommand{\nhforces}[2]{\nxforces{\hmodel}{#1}{#2}}

\newcommand{\wforces}[2]{\xforces{\wmodel}{#1}{#2}}
\newcommand{\nwforces}[2]{\nxforces{\wmodel}{#1}{#2}}



\newcommand{\sforce}[1]{\forcesym{#1}}
\newcommand{\snforce}[1]{\nforcesym{#1}}

\newcommand{\mforces}[2]{\forces{\model,#1}{#2}}
\newcommand{\muforces}[2]{\forces{\model,\layer,#1}{#2}}
\newcommand{\mwforces}[1]{\forces{\model,\world}{#1}}
\newcommand{\muwforces}[1]{\forces{\model,\layer,\world}{#1}}
\newcommand{\mnforces}[2]{\nforces{\model,#1}{#2}}
\newcommand{\munforces}[2]{\nforces{\model,\layer,#1}{#2}}
\newcommand{\mwnforces}[2]{\nforces{\model,\world,#1}{#2}}
\newcommand{\muwnforces}[1]{\nforces{\model,\layer,\world}{#1}}


\newcommand{\inrel}[3][\layer]{#2\Rel[#1]#3}

\newcommand{\defn}[1]{{\textit{\textbf{#1}}\/}}

%%%%%%%%%%%%%%%%%%%%%%%%%%%%%%%%%%%%%%%%%%%%%%%%%%%
%%TIKZ stuff

\tikzstyle{node}=[circle,draw]
\tikzstyle{s-node}=[circle,draw,inner sep=2pt]
\tikzstyle{nonode}=[inner sep=0pt]
\tikzstyle{access}=[-stealth,blue]
\tikzstyle{future}=[dashed,-stealth,thick]
%%\tikzstyle{suricata}=[future,red]
%%\tikzstyle{deleted}=[->,gray]
%%\tikzstyle{loop}=[->,red]

\newcommand{\forcebox}[2]{\parbox{1em}{\hbox to1em{\strut$#1$\hss} \hbox to1em{\strut$#2$\hss}}}
\newcommand{\mforcebox}[2]{\forcebox{\sforce{#1}}{{\color{red!70}{\snforce{#2}}}}}
\newcommand{\oforcebox}[2]{\forcebox{\sforce{#1}}{}}
\newcommand{\nforcebox}[2]{\forcebox{}{{\color{red!70}{\snforce{#2}}}}}

%%%%%%%%%%%%%%%%%%%%%%%%%%%%%%%%%%%%%%%%%%%%%%%%%%%%%

\renewcommand{\PROP}{\mathcal{V}}
%\newcommand{\FORM}{\mathcal{F}}

\newcommand{\formula}{\varphi}
\newcommand{\formulb}{\psi}

\newcommand{\IML}{\iml}
\newcommand{\IK}{\mathsf{IK}}
\newcommand{\IT}{\mathsf{IT}}


\newcommand{\kripke}{\mathcal{K}_\layer}
\newcommand{\kripkea}{\mathcal{K}_\layera}


\endinput
